\section{คู่มือสำหรับผู้พิจารณาคำร้อง}
\subsection{ขั้นตอนร่วม}
เปิด \textbf{การอนุมัติ/รับทราบ} แล้วเลือกบทบาทของตน badge แสดงจำนวนรายการรอ สามารถกรองตามสถานะ ประเภทลา หรือจำนวนวัน จากนั้นตรวจนิสิต เหตุผล ช่วงวัน หลักฐาน รายวิชา และประวัติขั้นตอน ปุ่มดำเนินการจะแสดงเมื่อคำร้องเป็น \code{CURRENT} ของบัญชีและ role นั้น ควรใส่ความคิดเห็น โดยเฉพาะการไม่อนุมัติ และตรวจกล่องยืนยัน

\subsection{อาจารย์ที่ปรึกษา}
ผู้มี \role{ADVISOR} เลือก \textbf{อาจารย์ที่ปรึกษา} ตรวจความสมเหตุสมผลและผลกระทบต่อการเรียน เมื่ออนุมัติจะส่งต่อหัวหน้าภาค เมื่อไม่อนุมัติ workflow สิ้นสุด เมนู \textbf{นิสิตในที่ปรึกษา} ใช้ตรวจรายชื่อ และ \textbf{ประวัติอนุมัติ} ใช้ดูรายการเดิม

\subsection{หัวหน้าภาควิชา}
ผู้มี \role{DEPARTMENT\_HEAD} เลือก \textbf{หัวหน้าภาควิชา} ตรวจคำร้อง หลักฐาน และผลจากที่ปรึกษา หากอนุมัติคำร้องไม่เกิน 15 วันจะเสร็จสมบูรณ์ หากเกิน 15 วันจะส่งต่อรองคณบดี

\subsection{รองคณบดี}
ผู้มี \role{ASSOCIATE\_DEAN} เลือก \textbf{รองคณบดีฯ} ซึ่งตาม logic ปัจจุบันจะแสดงคำร้องเกิน 15 วัน ตรวจข้อมูลและผลก่อนหน้า เมื่ออนุมัติคำร้องเป็น \code{APPROVED} เมื่อไม่อนุมัติเป็น \code{REJECTED}

\subsection{ปุ่มดำเนินการไม่ปรากฏ}
ตรวจว่าคำร้องถึงขั้นนี้หรือไม่ ถูกดำเนินการแล้วหรือไม่ บัญชีมี role ถูกต้องหรือไม่ และผู้อนุมัติใน snapshot ตรงกับบัญชีหรือไม่ หากยังผิดปกติให้แจ้งเลขคำร้องและ role แก่เจ้าหน้าที่
